\section{Implementación de un PID en un microcontrolador}
\label{sec:p3}
\enunciado{¿Cómo se implementa un sistema PID en un microcontrolador? ¿Por qué no se debe usar la
función delay como en Arduino? Muestre el código en un microcontrolador sin la función
delay.}

\subsection{De la ley continua a la ecuación en diferencias}
El PID continuo es
$u(t)=K_p e(t) + K_i\int_0^t e\,d\tau + K_d\,\dot e(t)$~\cite{astrom2006advanced}.
En el microcontrolador se evalúa una vez por periodo y cada término se discretiza por
separado. Se usa Euler hacia atrás, $s\approx(1-z^{-1})/T_s$, que es simple y estable para
cualquier $T_s$~\cite{astrom1997computer,ogata1996discreto}:
\begin{align}
  P[k] &= K_p\,e[k], \\
  I[k] &= I[k-1] + K_i T_s\,e[k], \label{eq:integral}\\
  D[k] &= \frac{T_f}{T_f+T_s}\,D[k-1] + \frac{K_d}{T_f+T_s}\big(\omega[k]-\omega[k-1]\big), \label{eq:derivada}\\
  u[k] &= \operatorname{sat}_{[u_\text{mín},\,u_\text{máx}]}\big(P[k]+I[k-1]+D[k]\big).
\end{align}
Cuando no hay saturación, el controlador equivale a
\begin{equation}
  C(z) = K_p + \frac{K_i T_s}{1-z^{-1}} + \frac{K_d\,(1-z^{-1})}{(T_f+T_s) - T_f z^{-1}}.
\end{equation}
El proyecto exige además las siguientes decisiones de implementación.
\begin{itemize}
  \item La derivada pura amplifica el ruido
  del encoder, por eso se aplica $K_d s/(T_f s+1)$ con $T_f\approx 2T_s$. Derivar $\omega$ en
  lugar de $e$ evita saltos cuando cambia la referencia~\cite{astrom2006advanced}.
  \item El mando se satura a los límites calibrados del freno, que van de la goma sin contacto al
  frenado máximo seguro.
  \item Como anti-windup se usa integración condicional (\emph{clamping}): si $u$ está en un límite
  y el error empuja hacia ese mismo límite, \eqref{eq:integral} no se actualiza. Con un freno
  unilateral esto es importante, porque cuando la referencia pide más frenado del posible, el integrador
  acumularía un mando que después mantiene el freno apretado aunque el disco ya esté por debajo
  de la curva, y el freno no puede devolver esa velocidad~\cite{astrom2006advanced,astrom2021feedback}.
  \item Los estados $I$, $D$ y $\omega[k-1]$ se borran al inicio de cada
  frenado para que los ensayos sean comparables.
\end{itemize}

\paragraph{Ganancias provisionales.}
Partiendo de la planta~\eqref{eq:planta_G} de la pregunta~5: como el frenado dura mucho menos
que $J_\text{eq}/b$, el polo en $-b/J_\text{eq}$ queda cerca del origen y, agregando el servo,
la planta del mando a la velocidad es aproximadamente un integrador con el polo del servo,
$G(s)\approx -a_\text{máx}/\big(s(\tau_\text{servo}s+1)\big)$, con $a_\text{máx}=K_f/J_\text{eq}$
por unidad de mando. El signo
negativo se compensa con la convención $e=\omega-\omega_\text{ref}$. Con
$a_\text{máx}\approx\SI{100}{rad/s^2}$, $K_p=0{,}04$ da un cruce
$\omega_c\approx K_p a_\text{máx}=\SI{4}{rad/s}$, el usado en la pregunta~2. El cero del PI se
coloca debajo del cruce, en $K_i/K_p=\SI{2}{rad/s}$ ($K_i=0{,}08$), y $K_d=5\times10^{-4}$ es
pequeño. Como el lazo resulta de tipo~2, sigue la rampa sin error en estado estacionario. Son
valores de partida que se rediseñarán con el modelo identificado.

\subsection{Estructura temporal: periodo fijo sin bloqueo}
El periodo $T_s$ usado para calcular $K_iT_s$, $K_d/T_s$ y los polos en $z$ tiene que ser el
periodo real. La implementación (Listados~\ref{lst:pid} y~\ref{lst:main}) separa tres
contextos~\cite{picosdk}:
\begin{enumerate}
  \item La interrupción del encoder (\texttt{gpio\_set\_irq\_enabled\_with\_callback}) se
  activa en ambos flancos y solo guarda \texttt{time\_us\_64()}, así que no se pierde ninguna marca.
  \item El temporizador periódico (\texttt{add\_repeating\_timer\_us(-20000, ...)}) usa un
  retardo negativo, que hace que el SDK programe cada llamado 20~ms después del \emph{inicio} del
  anterior, de modo que el tiempo de ejecución no se suma al periodo. En cada tic se mide
  $\omega$, se calcula $\omega_\text{ref}(kT_s)$, se ejecuta \texttt{pid\_paso}, se escribe el
  PWM del servo y se encola la muestra. Primero se lee y se actúa, y después se guarda el
  registro, para minimizar la latencia entre medición y mando.
  \item El bucle principal no tiene requisitos de tiempo. Lee la consola con
  \texttt{getchar\_timeout\_us(0)}, que no bloquea, y vacía la cola de telemetría por USB. Si la
  cola se llena se descarta telemetría y se cuenta, pero nunca se retrasa el control.
\end{enumerate}
El tiempo de ejecución de cada paso se registra (\texttt{ejecucion\_us}) para demostrar que es
mucho menor que $T_s$.

\subsection{¿Por qué no usar \texttt{delay()}?}
El patrón típico de Arduino \texttt{loop()\{ medir(); pid(); actuar(); Serial.println(...);
delay(20); \}} tiene tres problemas:
\begin{enumerate}
  \item El periodo real no es $T_s$. El ciclo dura $T_\text{delay}+T_\text{ejec}$, y
  $T_\text{ejec}$ varía con la ruta del código y, sobre todo, con la impresión serie: a
  115\,200~baudios, una línea CSV de unos 40 caracteres tarda $\approx\SI{3{,}5}{ms}$. Las
  ganancias discretas $K_iT_s$ y $K_d/T_s$ y la ubicación de los polos en $z$ quedan calculadas
  para un periodo que no es el real, y el \emph{jitter} añade ruido de muestreo no
  uniforme~\cite{astrom1997computer}.
  \item La base de tiempo deriva. Si la referencia se calcula como $\omega_0-a\,kT_s$
  con un contador de muestras, el error se acumula y la curva seguida resulta más lenta que
  la diseñada.
  \item Bloquea el procesador. Durante \texttt{delay()} el programa no atiende la
  consola, los botones, la telemetría ni una condición de seguridad, como cortar el freno si se
  pierde el encoder~\cite{arduinoblink}.
\end{enumerate}
Si se trabaja en Arduino, la alternativa es la temporización no bloqueante con
\texttt{millis()}/\texttt{micros()}~\cite{arduinoblink}. El instante siguiente se acumula
(\texttt{siguiente += TS}), en vez de tomarse desde el instante actual
(\texttt{siguiente = ahora + TS}), para no acumular deriva:
\begin{lstlisting}[caption={Patrón no bloqueante equivalente en Arduino.},label={lst:millis},numbers=none]
if ((int32_t)(micros() - siguiente) >= 0) {  // seguro ante desborde
    siguiente += TS_US;                       // sin deriva acumulada
    paso_de_control();                        // medir -> PID -> actuar
}
atender_serial_y_botones();                   // nunca se bloquea
\end{lstlisting}
En la Pico se prefiere el temporizador por hardware (Listado~\ref{lst:main}), porque además
garantiza el periodo aunque el bucle principal esté ocupado.

\lstinputlisting[caption={Módulo PID portable (\texttt{firmware/src/pid\_controller.c}), el mismo que usa la simulación.},label={lst:pid}]{\firmware/src/pid_controller.c}

\subsection{Evidencia \simulado}
Para comprobar el código sin el banco, \texttt{evidencias/simulacion\_pid.c} compila el mismo
\texttt{pid\_controller.c} contra una planta \emph{ilustrativa}:
$\dot\omega = -\omega/8 - 1{,}5 - 100\,b$, donde $b$ es la posición del servo, con
$\tau_\text{servo}=\SI{50}{ms}$ y una trama PWM de retardo. Integra a 0,1~ms y el control corre
a $T_s=\SI{20}{ms}$. Los parámetros son supuestos y no provienen de una identificación del banco. Las
figuras y métricas se regeneran con \texttt{./compilar.sh}.

\paragraph{Anti-windup.} La referencia pide \SI{140}{rad/s^2} durante 0,5~s, más de lo que el
freno puede dar ($\approx\SI{120}{rad/s^2}$), y luego \SI{25}{rad/s^2}
(Fig.~\ref{fig:aw}). Sin anti-windup el freno sigue apretado después de alcanzar la curva y el
disco cae cerca de \SI{20}{rad/s} por debajo; como el freno no puede acelerarlo, queda más de
1~s suelto ($u=0$) esperando a que la referencia lo alcance. Con clamping el IAE baja de \iaeAwNo{} a
\iaeAwSi{}~rad y la fracción de muestras saturadas pasa de \satAwNo{} a \satAwSi{}\,\%.

\begin{figure}[h]
\centering
\begin{tikzpicture}
\begin{groupplot}[group style={group size=1 by 2, vertical sep=0.6cm, x descriptions at=edge bottom},
  width=0.92\linewidth, height=4.6cm, xmin=0, xmax=4, grid=major, grid style={gray!20},
  legend style={font=\footnotesize, at={(0.99,0.97)}, anchor=north east}]
\nextgroupplot[ylabel={$\omega$ [rad/s]}]
  \addplot[black, dashed, thick] table[x=t, y=ref, col sep=comma] {evidencias/generado/aw_si.csv};
  \addlegendentry{$\omega_\text{ref}$}
  \addplot[blue!70!black, thick] table[x=t, y=omega, col sep=comma] {evidencias/generado/aw_si.csv};
  \addlegendentry{con anti-windup}
  \addplot[red!70!black, thick] table[x=t, y=omega, col sep=comma] {evidencias/generado/aw_no.csv};
  \addlegendentry{sin anti-windup}
\nextgroupplot[ylabel={$u$ [norm.]}, xlabel={$t$ [s]}, ymin=-0.05, ymax=1.1]
  \addplot[blue!70!black, thick, const plot] table[x=t, y=u, col sep=comma] {evidencias/generado/aw_si.csv};
  \addplot[red!70!black, thick, const plot] table[x=t, y=u, col sep=comma] {evidencias/generado/aw_no.csv};
\end{groupplot}
\end{tikzpicture}
\caption{\simulado{} Seguimiento con una referencia que satura el freno al inicio.}
\label{fig:aw}
\end{figure}

\paragraph{Temporizador frente a \texttt{delay()}.} Con la rampa alcanzable de
\SI{40}{rad/s^2}, el mismo PID se ejecuta con el temporizador ($T_s$ exacto) y con un bucle
\texttt{delay(20)} más un cálculo de 0,3~ms y una impresión serie de 2,5 a 4,5~ms
(Fig.~\ref{fig:delay}). El periodo real medio sube a \periodoDelay~ms (máximo
\periodoMaxDelay~ms). Como la referencia avanza con el contador de muestras, al final el
firmware está atrasado \desfaseDelay~ms respecto al reloj real y la desaceleración efectiva
baja a $\approx$\decelEfectivaDelay~rad/s$^2$. Respecto de la curva diseñada, el IAE pasa de
\iaeTimer{} a \iaeDelay~rad y el RMSE de \rmseTimer{} a \rmseDelay~rad/s. Como las ganancias y
el código son idénticos, la diferencia se debe solo a la temporización.

\begin{figure}[h]
\centering
\begin{tikzpicture}
\begin{groupplot}[group style={group size=2 by 1, horizontal sep=1.6cm},
  height=5cm, grid=major, grid style={gray!20},
  legend style={font=\footnotesize, at={(0.99,0.97)}, anchor=north east}]
\nextgroupplot[width=0.56\linewidth, xmin=0, xmax=4.5, xlabel={$t$ [s]}, ylabel={$\omega$ [rad/s]}]
  \addplot[black, dashed, thick] table[x=t, y=ref, col sep=comma] {evidencias/generado/timer.csv};
  \addlegendentry{$\omega_\text{ref}(t)$}
  \addplot[blue!70!black, thick] table[x=t, y=omega, col sep=comma] {evidencias/generado/timer.csv};
  \addlegendentry{temporizador}
  \addplot[red!70!black, thick] table[x=t, y=omega, col sep=comma] {evidencias/generado/delay.csv};
  \addlegendentry{\texttt{delay(20)}}
\nextgroupplot[width=0.42\linewidth, xlabel={muestra $k$}, ylabel={periodo real [ms]}, ymin=19, ymax=26]
  \addplot[blue!70!black, thick, domain=0:200] {20};
  \addplot[red!70!black, only marks, mark size=0.6pt] table[x=k, y=periodo_ms, col sep=comma] {evidencias/generado/periodos_delay.csv};
\end{groupplot}
\end{tikzpicture}
\caption{\simulado{} Izquierda: seguimiento de la rampa. Derecha: periodo entre muestras con \texttt{delay(20)} (rojo) frente al temporizador (azul).}
\label{fig:delay}
\end{figure}

\pendiente{Repetir ambas comparaciones en el banco cuando la planta esté armada (preguntas 6 y
7): registrar \texttt{ejecucion\_us}, el periodo real entre muestras y el seguimiento con la
planta física.}
